%%%% INSERT_AFTER_LINE: 76
%%%% ANCHOR_TEXT: % ---- END NOTES ----
% CODE INPUT
The first half of this thesis measures how QSAR models behave as label noise rises (Chapter~\ref{chap:robustness}). We define robustness as the normalised area under the curve of $R^2$ retained as noise increases. Under Gaussian noise on QM9, model architecture explains 49.5\% of the variance in robustness, and representation explains 9.2\%. For $R^2$ at noise level 1.0, where the noise has the same spread as the clean training labels, the two shares are comparable: 29.1\% for architecture and 26.6\% for representation. These shares are measured under noise. \textbf{[RESULT NEEDED: the same variance split on $R^2$ before any noise is added.]} Systematic noise did the most damage: noise shared within a scaffold family, and labels clipped at an assay limit. The shape of the noise distribution barely mattered. Under Gaussian noise on QM9, mean predicted uncertainty rose with noise for every probabilistic model and representation. Per-molecule uncertainty, however, did not reliably pick out which labels were corrupted (Section~\ref{sec:uq}). NGBoost and the Gaussian process fit a separate scale for each molecule. They gave the strongest per-molecule correlation between predicted uncertainty and injected noise. Under censoring, models were most confident on the labels that had been clipped. Uncertainty alone therefore cannot tell us which labels to fix. This led us to methods that judge each label in another way.

% CODE INPUT
Learning with noisy labels (LNL) is the broad family of methods for training when some labels are wrong. Noise mitigation, which removes, reweights or corrects labels, is one part of it. In an exploratory study, we adapted twelve LNL methods to regression and tested them in 221 configurations (Section~\ref{sec:lnl_motivation}). The task was the QM9 HOMO--LUMO gap under Gaussian noise. Several variants judged a label by the labels of its neighbours, found with a distance between molecules. All 20 configurations with the lowest retention of $R^2$ under noise used such a distance. This held for every distance we tried: Tanimoto on ECFP4, Euclidean on physicochemical descriptors, and Mahalanobis. The methods that held up estimated noise from the model's own loss or uncertainty, not from neighbours. The study covers one dataset and one noise type, so we treat it as motivation rather than as a general result.

% CODE INPUT
We hypothesised that the distance itself was at fault. On a given dataset, most features of a general-purpose representation either do not vary or do not relate to the property. Every such feature adds a term to the distance, so together these irrelevant features dominate it. Their contribution grows linearly with the number of features, not exponentially. Chapter~\ref{chap:flaws} measures this on base representations and on the distances computed from them. The fix we propose is to select a fixed number of task-relevant features from several representations, and to compute distances in that compact space. A first version showed that these selected feature sets also work as ordinary model inputs, with no distance involved. The question then became how to choose the features, which is the subject of Chapter~\ref{chap:method}. Chapter~\ref{chap:similarity} returns to distances, noise detection and LNL, including noise mitigation, with the hybrid representation as the feature space.
%%%% END

%%%% INSERT_AFTER_LINE: 141
%%%% ANCHOR_TEXT: % ---- END NOTES ----
% CODE INPUT
Chapter~\ref{chap:background} introduces, once, the ideas that the later chapters rely on: molecular representations, model families, molecular similarity and activity cliffs, noise in molecular data, and the evaluation protocol. It defines the seven noise conditions used in Chapter~\ref{chap:robustness}. It ends by introducing NoiseInject, the open-source benchmarking framework released with this work.

% CODE INPUT
Chapter~\ref{chap:robustness} is adapted from the paper ``NoiseInject: A Benchmarking Framework to Study Noise Robustness in QSAR Models'' by Punt, Hanser, Werner and Morris, in preparation for the \textit{Journal of Cheminformatics}. The main task is the HOMO--LUMO gap on QM9. Its labels are computed, so they carry negligible measurement noise. We repeat the experiments on three assay datasets with experimental labels: LogD, Caco-2 and hERG $K_i$. The chapter splits the variance in performance, and in how much performance is retained under noise, between representation and model architecture. It compares seven noise conditions at a matched amount of noise. It also compares each model with its probabilistic counterpart, and asks what label noise does to predicted uncertainty. Model architecture matters more than representation for noise robustness. For $R^2$ under noise, the two matter about equally. The shape of the noise distribution makes little difference. Noise shared within a scaffold family, and labels clipped at an assay limit (censoring), do the most damage. NGBoost and random forests were the most robust models on QM9, and random forests were the most consistent across all four datasets. Predicted uncertainty rises with noise on average, but it does not pick out which labels were corrupted.

% CODE INPUT
Chapter~\ref{chap:flaws} is a motivation chapter. It introduces no method, and its measurements are smaller than those of Chapter~\ref{chap:robustness}. It makes three arguments. First, on a given dataset, standard representations carry many features that do not vary, and some map different molecules to the same vector. Second, similarity computed over those representations inherits both problems. Third, methods that rely on that similarity fail as a result. One such method cleans labels by comparing each molecule's label with the labels of its nearest neighbours. It is a noise mitigation method tested in the exploratory study of learning with noisy labels, and that study comes from my confirmation report. The chapter ends by stating what a representation would need in order to fix these problems. That is what Chapter~\ref{chap:method} builds.

% CODE INPUT
Chapter~\ref{chap:method} is planned as the first of two method papers, intended for submission to a venue still to be confirmed. From several base representations, the method selects each one's most important features for the task. It splits a fixed budget of features between the representations and combines the selected features into one representation. The chapter also describes the datasets and the experimental setup. It will test when the hybrid helps, against the best single representation and against an ensemble of one model per representation. It will test why it helps, against PCA, mRMR, recursive feature elimination and random features at the same number of features. We will also ask whether the selected features can be read against the chemistry, and how stable the selection is across folds. A further test asks whether search settings found on one dataset carry over to another. The chapter will test whether the method extends beyond small molecules. Antibodies are covered in most depth, with reactions, nucleotide sequences and cancer genomics as proofs of concept. Finally, it will test whether the method extends to data held at separate sites that cannot be pooled. \textbf{[RESULT NEEDED: one sentence of findings for Chapter 5.]}

% CODE INPUT
Chapter~\ref{chap:similarity} is planned as the second method paper, also intended for submission to a venue still to be confirmed. It returns to the failures described in Chapter~\ref{chap:flaws} and asks whether a hybrid representation repairs them. For similarity, it compares standard representations with hybrids selected for predicting each molecule's activity from its five nearest neighbours. For noise detection and estimation, it asks whether neighbours in a hybrid pick out corrupted labels better than neighbours in a standard representation. It then tests learning with noisy labels, including noise mitigation, one noise type at a time. \textbf{[RESULT NEEDED: one sentence of findings for Chapter 6.]}

% CODE INPUT
Chapter~\ref{chap:discussion} summarises each chapter. Its closing argument is that the noise should be modelled before a method is chosen. No single method fits every dataset, and simple methods are often as good as complex ones. The chapter then sets out the limitations of this work and directions for future work.

% CODE INPUT
This work contributes publications and software. The publications are the paper behind Chapter~\ref{chap:robustness}, now in preparation, and the planned papers from Chapters~\ref{chap:method} and~\ref{chap:similarity}. They are listed under Research Papers in the Summary of Research Contributions. The software is NoiseInject, which is open source, and the KIRBy software repository, which implements the hybrid search and its experiments. Repositories, licences and data are described in Appendix~\ref{app:availability}.
%%%% END

%%%% INSERT_AFTER_LINE: 145
%%%% ANCHOR_TEXT: % [notes moved to thesis_notes_moved.md]
% CODE INPUT
This DPhil was funded by the University of Oxford's UK EPSRC Centre for Doctoral Training in Sustainable Approaches to Biomedical Science: Responsible and Reproducible Research, grant EP/S024093/1. It was also supported by Lhasa Limited, a not-for-profit organisation.

% CODE INPUT
Thierry Hanser and Stephane Werner at Lhasa Limited co-supervised this work, with Garrett Morris as supervisor. Chapter~\ref{chap:robustness} is adapted from the jointly authored NoiseInject manuscript (Punt, Hanser, Werner and Morris). Each author's contribution is set out in the Statement of Originality. Lhasa runs Effiris, a federated-learning platform \citep{Bassani2023}. Lhasa also supplied the hERG data used in Chapter~\ref{chap:method}: 7,772 training molecules and 4,383 test molecules. We use these as data held at separate sites, the setting that federated learning addresses \citep{Bassani2023, Hanser2023}.

% CODE INPUT
The Oxford Protein Informatics Group (OPIG) pointed us to the antibody and protein datasets used in Chapters~\ref{chap:method} and~\ref{chap:similarity}. The full list of datasets is in Appendix~\ref{app:datasets}. The collaboration also produced five posts on OPIG's blog, Blopig, listed under Blogposts in the Summary of Research Contributions.

% CODE INPUT
The adaptive hybrid representation, its software and the experiments in Chapters~\ref{chap:flaws} to~\ref{chap:similarity} are my own work.
%%%% END
